\section{Project Management}
\label{sec:mgmt}

\subsection{Team and Responsibilities}
The five team members split into sub-teams by subsystem, and three additional
contributors assisted with the radio and sensors.
Table~\ref{tab:team} lists each member's responsibilities, taken from the
proposal schedule and the fall firmware task plan. The proposal organized the work into four sub-teams:
\begin{itemize}
\item \emph{LoRa}: Harrison LeTourneau, Jared Pratt, Adam Welsh, and Aliou
Tippett designed, implemented, and tested the radio link between the vessel
and the ground station.
\item \emph{Sensor}: Corbin Barney, Charbel Salloum, Connor Stevens, and
Derek Tinoco selected, integrated, and tested the GPS, IMU, and wind vane.
\item \emph{Embedded}: all five members, with Derek Tinoco, Adam Welsh, and
Aliou Tippett, built the bench test setups for the radio and sensor suite.
\item \emph{ALAN}: the five members shared vessel integration, mechanical
assembly, firmware integration, and on-water testing.
\end{itemize}
Adam Welsh and Aliou Tippett are students at the University of Utah, and
Derek Tinoco is a computer scientist who assisted with early sensor
integration.

\begin{table}[!htb]
\caption{Team Responsibilities}
\label{tab:team}
\centering
\footnotesize
\begin{tabular}{@{}p{0.22\columnwidth}p{0.3\columnwidth}p{0.38\columnwidth}@{}}
\toprule
Member & Mechanical & Firmware and software \\
\midrule
Charbel Salloum & Wingsail & Wingsail control, live tack planner, path planning \\
Connor Stevens & Wingsail & Wingsail control, telemetry \\
Corbin Barney & Keel, rudder & Sensor drivers, shared data types, positioning, emergency recovery \\
Harrison LeTourneau & Hull & LoRa link (CM4), commander, positioning, ground station \\
Jared Pratt & Keel, rudder & Heading and sail PIDs, arbiters, sail planner, ground station, path planning \\
\bottomrule
\end{tabular}
\end{table}

\subsection{Schedule}
Fig.~\ref{fig:gantt} shows the schedule planned in the proposal. The spring
and summer covered fabrication and assembly, and the fall covered firmware,
control, navigation, and water testing. Table~\ref{tab:schedule} gives the
owner assigned in the proposal and the status of each milestone at the time
of writing.

\subsection{Cost}
The proposal estimated a total of about \$1000, shown against actual
spending in Table~\ref{tab:cost}. The largest single expense was expected to be
epoxy resin, and the most expensive component the hull, because of the
printed mold and its filament. On-campus groups sponsored some materials,
including 3D printing filament, and university staff loaned reusable
equipment such as respirators. John Davies provided a space in the ECE Senior Capstone Lab to house materials and a space to assemble the project.

\begin{table}[!htb]
\caption{Estimated and Actual Cost}
\label{tab:cost}
\centering
\begin{tabular}{@{}lrr@{}}
\toprule
Category & Estimate (\$) & Actual (\$) \\
\midrule
Hull materials (resin, fiberglass, mold) & 400 & 375 \\
Keel, rudder, sail/tail (steel, carbon fiber, etc.)    & 200 & 145 \\
Electronics                              & 200 & 250 \\
Contingency                              & 200 & 250 \\
\midrule
Total                                    & 1000 &  1020\\
\bottomrule
\end{tabular}
\end{table}